%======================================================================
%  Galamsey — Academic Beamer slides (Overleaf-ready, self-contained)
%
%  Mirrors the content of galamsey_motivation.qmd, minus dynamic
%  elements (the animated GIF is replaced by its final cumulative frame).
%  Regression results are shown as proper LaTeX tables.
%
%  Overleaf setup — upload alongside this .tex a "figures/" folder that
%  mirrors the repo's outputs/figures/ tree, i.e.:
%
%    galamsey_motivation_slides.tex
%    figures/
%      d1c_districts_ts_top12.png  d1d_districts_extent_top20.png
%      d1d_national_gold_price.png
%      maps/hex5km/      d1a_hex5km_cumulative.png  d1b_hex5km_artisanal_vs_industrial.png
%      maps/districts/   d1a_districts_cumulative.png  d1b_districts_artisanal_vs_industrial.png
%      spatial_clustering/  d2_map_incidence.png  d2a_morans_by_year.png  d2a_lisa_2017.png
%                           d2bc_null_histogram.png  d2e_schematic.png  d2f_event_study.png
%                           d2_geography_firststage.tex  d2d_spatial_lag.tex   (regression tables)
%      motivating_facts/ d5a_ndvi_evi_distance_binscatter.png  d5b_ndvi_evi_event_study.png
%
%  Images are referenced by basename (resolved via \graphicspath below);
%  the two regression-table fragments are pulled in with \input using their
%  figures/spatial_clustering/ path. Both must live under the project root —
%  Overleaf cannot read paths above it (no "../").
%
%  Theme: built-in Beamer base recoloured to the galamsey wine palette
%  (#450457) — compiles on Overleaf and any standard TeX install with no
%  extra theme packages. Remove "handout" below to enable \pause overlays.
%======================================================================
\documentclass[aspectratio=169,handout]{beamer}

\usepackage[utf8]{inputenc}
\usepackage[T1]{fontenc}
\usepackage{lmodern}
\usepackage{amsmath,mathtools}
\usepackage{booktabs}
\usepackage{graphicx}
\usepackage{xcolor}
\usepackage{tikz}
\usepackage{caption}
\captionsetup{font=scriptsize,labelfont=bf}

\graphicspath{%
  {figures/}%
  {figures/maps/hex5km/}%
  {figures/maps/hex10km/}%
  {figures/maps/districts/}%
  {figures/spatial_clustering/}%
  {figures/motivating_facts/}%
}

%---------------------------------------------------------------
%  Palette — galamsey wine
%---------------------------------------------------------------
\definecolor{galawine}{RGB}{69,4,87}
\definecolor{galaorange}{RGB}{230,126,34}
\definecolor{galagrey}{RGB}{120,120,120}

\usefonttheme{default}
\setbeamertemplate{navigation symbols}{}
\setbeamercolor{structure}{fg=galawine}
\setbeamercolor{title}{fg=galawine}
\setbeamercolor{frametitle}{bg=galawine,fg=white}
\setbeamercolor{alerted text}{fg=galaorange}
\setbeamercolor{block title}{bg=galawine,fg=white}
\setbeamercolor{block body}{bg=galawine!6,fg=black}
\setbeamercolor{itemize item}{fg=galawine}
\setbeamercolor{itemize subitem}{fg=galaorange}
\setbeamertemplate{itemize items}[circle]
\setbeamersize{text margin left=0.7cm,text margin right=0.7cm}

% Flat wine frametitle bar (metropolis-like)
\setbeamertemplate{frametitle}{%
  \nointerlineskip%
  \begin{beamercolorbox}[wd=\paperwidth,ht=3.0ex,dp=1.4ex,leftskip=0.5cm,rightskip=0.5cm]{frametitle}%
    \usebeamerfont{frametitle}\insertframetitle%
  \end{beamercolorbox}%
}

% Full-bleed wine section divider frames (echo the QMD purple "Part" slides)
\AtBeginSection[]{%
  \begin{frame}[plain]
    \begin{tikzpicture}[remember picture,overlay]
      \fill[galawine] (current page.south west) rectangle (current page.north east);
    \end{tikzpicture}
    \vfill
    \centering{\color{white}\usebeamerfont{frametitle}\Large\textbf{\insertsection}}
    \vfill
  \end{frame}
}

\newcommand{\srcnote}[1]{{\scriptsize\color{galagrey}#1}}

%---------------------------------------------------------------
%  Title
%---------------------------------------------------------------
\title{Galamsey}
\subtitle{Spatial Distribution, Clustering, and Evidence for Economic Spillovers}
\author{Sree Ayyar \and Hema Balarama \and Yoshiki Wiskamp \and Till Meissner}
\institute{London School of Economics}
\date{\today}

\begin{document}

\maketitle

%======================================================================
\begin{frame}{Overview}
  \textbf{Data:} Barenblitt et al.\ (2021) remote-sensing panel --- mine
  extent polygons, annual 2007--2017, SW Ghana ($\sim$104{,}730 km$^2$).

  \medskip
  \textbf{Unit of analysis:} 5 km hexagonal grid over the study region.

  \medskip
  \textbf{Questions:}
  \begin{enumerate}\setlength\itemsep{0.4em}
    \item \textbf{Where?} Spatial distribution of galamsey across SW Ghana
    \item \textbf{Clustered?} Statistical evidence for spatial autocorrelation
    \item \textbf{Why clustered?} Geography vs.\ social and economic drivers
    \item \textbf{Spreading?} Evidence for spillovers along river corridors
  \end{enumerate}

  \medskip
  \alert{Core question:} Is the spatial clustering of galamsey explained by
  underlying geography (gold geology, river networks), or is there evidence
  of social and economic spillover processes?
\end{frame}

%======================================================================
\section{Where does galamsey happen?}
%======================================================================

\begin{frame}{Data: Barenblitt et al.\ (2021)}
  \begin{columns}[T]
    \begin{column}{0.56\textwidth}
      \textbf{Remote-sensing derived mine extent, SW Ghana}
      \begin{itemize}\setlength\itemsep{0.45em}
        \item \textbf{Annual time series:} 2007--2017, new conversions per year
        \item \textbf{Full 2019 extent:} artisanal vs industrial mining
        \item Classified from Landsat imagery, Random Forest model
        \item \textbf{Coverage:} SW Ghana only ($\sim$104{,}730 km$^2$) ---
              northern districts have no data, not zero mines
        \item \textbf{Accuracy:} producer accuracy 75.6\% --- $\sim$1 in 4
              real mines missed; treat as positive labels only
      \end{itemize}
    \end{column}
    \begin{column}{0.44\textwidth}
      \centering
      \includegraphics[width=\linewidth]{d2_map_incidence.png}
    \end{column}
  \end{columns}
\end{frame}

\begin{frame}{Fine-Scale Incidence --- 5 km Hexagons}
  \centering
  \includegraphics[height=0.74\textheight]{d1a_hex5km_cumulative.png}

  \srcnote{Cumulative artisanal mining land cover per 5 km hexagon, 2007--2017
  (final frame of the animated sequence in the HTML deck). Tight spatial
  clustering is evident within and across district boundaries.}
\end{frame}

\begin{frame}{Artisanal vs Industrial Mining --- 5 km Hexagons}
  \centering
  \includegraphics[width=0.92\linewidth]{d1b_hex5km_artisanal_vs_industrial.png}

  \srcnote{Left: artisanal (galamsey). Right: industrial (formal). Artisanal
  mining is more dispersed and river-proximate; industrial concentrates in
  fewer, larger sites. Barenblitt 2019 extent.}
\end{frame}

\begin{frame}{District-Level Extent}
  \begin{columns}[T]
    \begin{column}{0.52\textwidth}
      \textbf{Galamsey is widespread across SW Ghana}
      \begin{itemize}\setlength\itemsep{0.45em}
        \item \textbf{96 of Ghana's 260 districts} fall within the Barenblitt
              survey area
        \item Artisanal mining is recorded across the full survey area ---
              not confined to a handful of hotspots
        \item The remaining 164 districts have no coverage: absent data means
              \emph{no survey}, not zero mines
        \item \textbf{Highly concentrated:} the top 10\% of districts
              ($\sim$10) account for \textbf{86\%} of total artisanal extent
      \end{itemize}
    \end{column}
    \begin{column}{0.48\textwidth}
      \centering
      \includegraphics[width=0.92\linewidth]{d1a_districts_cumulative.png}
    \end{column}
  \end{columns}

  \medskip
  \srcnote{Cumulative artisanal mining land cover per district, 2005--2019.
  Plasma colour scale (sqrt transform). Barenblitt et al.\ (2021).}
\end{frame}

%======================================================================
\section{Is galamsey spatially clustered?}
%======================================================================

\begin{frame}{Global Moran's I: Is Clustering Statistically Significant?}
  \centering
  \includegraphics[height=0.72\textheight]{d2a_morans_by_year.png}

  \srcnote{Global Moran's I on annual new mining (ha) per 5 km hexagon,
  2007--2017. Queen contiguity weights, row-standardised. Red = $p<0.05$.
  Moran's I consistently positive and significant.}
\end{frame}

\begin{frame}{LISA Clusters: Where Are the Hotspots?}
  \centering
  \includegraphics[height=0.72\textheight]{d2a_lisa_2017.png}

  \srcnote{Local Moran's I clusters (LISA), artisanal mining extent 2019.
  \textbf{HH} = high mining surrounded by high mining (hotspot);
  \textbf{LL} = low surrounded by low. HH clusters are tightly localised
  in a corridor of SW Ghana.}
\end{frame}

%======================================================================
\section{What drives the clustering?}
%======================================================================

\begin{frame}{Decomposing Spatial Clustering}
  The observed spatial autocorrelation has two components:
  \[
    I_{\text{observed}} \;=\;
    \underbrace{I_{\text{geography}}}_{\text{clustering due to geography}}
    \;+\;
    \underbrace{I_{\text{excess}}}_{\text{excess (unexplained)}}
  \]
  \begin{itemize}\setlength\itemsep{0.5em}
    \item \alert{$I_{\text{geography}}$} --- mines locate where geology and
          rivers make mining viable: spatial concentration, but not
          necessarily neighbour-to-neighbour spillover.
    \item \alert{$I_{\text{excess}}$} --- mines cluster \emph{more than
          geography alone can explain}. Consistent with social/economic
          spillovers: labour mobility, information diffusion, reduced
          enforcement risk.
    \item \textbf{Strategy:} (1) confirm geography is a genuine predictor;
          (2) show residual clustering far exceeds what the richest geography
          null --- including an unrestricted spatial spline (M5) --- can
          account for.
  \end{itemize}
\end{frame}

\begin{frame}{Geography as a Predictor: First Stage}
  \begin{center}
    {\footnotesize\input{figures/spatial_clustering/d2_geography_firststage.tex}}
  \end{center}
  \srcnote{LPM, 5 km hexes, robust SEs. $^{***}\,p<0.01$.}
  \begin{itemize}\setlength\itemsep{0.35em}
    \item \textbf{River proximity:} mining concentrates near waterways ---
          alluvial gold deposits and water requirements of artisanal processing.
    \item \textbf{Gold geology:} hexes on gold-suitable ground show higher
          mine incidence --- consistent with ore-targeting.
    \item Low $R^2$: geography explains \emph{where}, not \emph{why} mines cluster.
  \end{itemize}
\end{frame}

\begin{frame}{Building the Geography Control: M1 $\to$ M4}
  \begin{itemize}\setlength\itemsep{0.5em}
    \item \textbf{M1} --- first-stage geography: distance to river + gold
          suitability
    \item \textbf{M2} --- adds neighbourhood geography (spatial lags of both
          covariates)
    \item \textbf{M3} --- adds terrain (elevation, slope)
    \item \textbf{M4} --- adds interactions
  \end{itemize}
  \medskip
  \begin{center}
    \IfFileExists{figures/spatial_clustering/d2_m1m4_ladder.tex}{%
      {\small\input{figures/spatial_clustering/d2_m1m4_ladder.tex}}%
    }{%
      \textit{[M1--M4 fit table pending --- run \texttt{d03c\_maup\_robustness.R}]}%
    }
  \end{center}
  \srcnote{5 km hexes. Fit rises M1$\to$M4 yet $p$\,(excess) $\approx 0$
  throughout. Robust at 1 km and 2 km (MAUP).}
\end{frame}

\begin{frame}{Ruling Out Unobserved Geography: The Spatial Spline (M5)}
  \begin{columns}[T]
    \begin{column}{0.52\textwidth}
      \[
        \text{logit}\,P(\text{mine}_i) =
          \underbrace{X_i\beta}_{\text{M4 cov.}} +
          \underbrace{f(\text{east}_i,\;\text{north}_i)}_{\text{smooth surface}}
      \]
      \begin{itemize}\setlength\itemsep{0.4em}
        \item A spatially \textbf{smooth} omitted driver --- geographic,
              geological, economic, or institutional --- could still explain
              the clustering even after M1--M4.
        \item \textbf{M5} adds a \textbf{thin-plate spline}: a flexible 2-D
              surface absorbing \emph{all} smooth spatial variation, measured
              or not.
        \item Lets \emph{pure location} soak up whatever covariates miss ---
              the most conservative geography null possible.
      \end{itemize}
    \end{column}
    \begin{column}{0.48\textwidth}
      \begin{center}
        \IfFileExists{figures/spatial_clustering/d2_m5_stats.tex}{%
          {\small\input{figures/spatial_clustering/d2_m5_stats.tex}}%
        }{%
          \textit{[M4 vs M5 stats table pending]}%
        }
      \end{center}
      \srcnote{5 km hexes. $p$\,(excess) $= 0$ under both M4 and M5.
      Robust at 2 km.}
    \end{column}
  \end{columns}
\end{frame}

\begin{frame}{M5: Observed Clustering Still Exceeds the Null}
  {\small\textbf{Simulation:} draw mine placements from model-predicted
  probabilities (no spillover term), repeat 500$\times$, compute Moran's I ---
  this is where clustering lands if only geography drives it.}

  \medskip
  \begin{columns}[T]
    \begin{column}{0.60\textwidth}
      \begin{center}
        \IfFileExists{figures/spatial_clustering/d2_m5_null_5km.png}{%
          \includegraphics[width=\linewidth]{d2_m5_null_5km.png}%
        }{%
          \textit{[Null distribution figure pending --- run \texttt{d03c\_maup\_robustness.R}]}%
        }
      \end{center}
    \end{column}
    \begin{column}{0.40\textwidth}
      \textbf{Residual Moran's I}
      \begin{center}
        \IfFileExists{figures/spatial_clustering/d2_m5_resid_moran.tex}{%
          {\small\input{figures/spatial_clustering/d2_m5_resid_moran.tex}}%
        }{%
          \textit{[Table pending]}%
        }
      \end{center}
    \end{column}
  \end{columns}
\end{frame}

\begin{frame}{Mining Propensity Net of Measured Covariates}
  \centering
  \IfFileExists{figures/spatial_clustering/d2_m5_surface_5km.png}{%
    \includegraphics[height=0.68\textheight]{d2_m5_surface_5km.png}%
  }{%
    \textit{[M5 surface map pending --- run \texttt{d03c\_maup\_robustness.R}]}%
  }

  \srcnote{M5 spatial term $s$(easting, northing), 5 km hexagons, mean-centred
  logit. \textbf{Red} = more mining than geology/river/terrain predicts at that
  hex; \textbf{Blue} = less. Not attributable to any smooth confounder.}
\end{frame}

%======================================================================
\section{Does galamsey spread? Evidence for spillovers}
%======================================================================

\begin{frame}{Spatial Lag Model: Estimating Equation}
  \[
    \Delta y_{it} \;=\; \alpha_i \;+\; \gamma_t
      \;+\; \rho\,\bar{y}_{N(i),\,t-1} \;+\; \varepsilon_{it}
  \]
  \begin{itemize}\setlength\itemsep{0.4em}
    \item $\Delta y_{it}$: new artisanal mining in hex $i$, year $t$ (ha; or
          0/1 for the LPM variant)
    \item $\bar{y}_{N(i),\,t-1}$: lagged \emph{cumulative} mining stock in
          queen-contiguous neighbours at end of $t-1$
    \item $\rho$: spatial spillover coefficient
    \item $\alpha_i$: hex FE --- absorb all time-invariant geography
    \item $\gamma_t$: year FE --- absorb national shocks (gold price, policy)
    \item \textbf{Identification:} within-hex variation over time. Residual
          concern: time-varying \emph{local} confounders (e.g.\ new roads).
  \end{itemize}
\end{frame}

\begin{frame}{Spatial Lag Regression: Results}
  \begin{center}
    {\footnotesize\input{figures/spatial_clustering/d2d_spatial_lag.tex}}
  \end{center}
  \begin{itemize}\setlength\itemsep{0.35em}
    \item \textbf{Intensive margin:} $+10$ ha cumulative neighbour mining
          $\to +0.91$ ha new mining this year ($t=7.5$, $p<10^{-13}$)
    \item \textbf{Extensive margin:} $+50$ ha in neighbours $\to \sim$13 pp
          higher probability of any mine onset ($t=8.8$, $p<10^{-17}$)
  \end{itemize}
  \srcnote{\textbf{Notes:} \texttt{fixest::feols}, two-way FE, SEs clustered by
  hex. $^{***}\,p<0.01$. Caveat: time-varying local confounders not ruled out
  without an instrument.}
\end{frame}

\begin{frame}{Upstream vs Downstream: Conceptual Logic}
  \begin{columns}[T]
    \begin{column}{0.44\textwidth}
      \textbf{Mechanism:} mining upstream contaminates waterways and degrades
      agricultural land --- reducing farming returns downstream and pushing
      communities into artisanal mining.

      \medskip
      \textbf{Prediction:} if this channel operates, upstream neighbours'
      mining should \emph{precede} focal-hex onset; downstream onset follows
      once agricultural productivity falls.
    \end{column}
    \begin{column}{0.56\textwidth}
      \centering
      \includegraphics[width=0.85\linewidth]{d2e_schematic.png}
    \end{column}
  \end{columns}
\end{frame}

\begin{frame}{Upstream vs Downstream: Evidence}
  \centering
  \includegraphics[height=0.66\textheight]{d2f_event_study.png}

  \srcnote{Mean cumulative mining in upstream vs downstream neighbours,
  normalised to the year before focal-hex onset ($t=-1=0$). If upstream mining
  degrades agricultural land and pushes downstream communities into mining,
  upstream neighbours' mining should rise \emph{before} the focal onset.
  River-adjacent 5 km hexes. Barenblitt et al.\ (2021).}
\end{frame}

\begin{frame}{Summary: A Cumulative Case for Economic Spillovers}
  \begin{table}
    \centering\small
    \begin{tabular}{p{0.24\textwidth}p{0.36\textwidth}p{0.30\textwidth}}
      \toprule
      \textbf{Evidence} & \textbf{What it shows} & \textbf{Geography alternative} \\
      \midrule
      M1--M4 null & $p$\,(excess) $= 0$ across all measured-geography models
        & Finer-scale geology than Girard 1:10M \\[0.3em]
      M5 spatial spline & Excess clustering survives the most flexible
        geography null & Non-smooth sub-grid driver \\[0.3em]
      Spatial lag & Neighbour stock predicts onset, within-hex $+$ year FE
        & Time-varying local confounders \\[0.3em]
      Direction of spread & Raw data: faster downstream; statistical tests
        inconclusive & --- \\
      \bottomrule
    \end{tabular}
  \end{table}
\end{frame}

%======================================================================
\section{Motivating Facts}
%======================================================================

\begin{frame}{National Trend vs Gold Price}
  \centering
  \includegraphics[height=0.70\textheight]{d1d_national_gold_price.png}

  \srcnote{Annual new galamsey conversion (ha) alongside international gold
  price (USD/troy oz). The 2007--2012 price boom coincides with a marked
  acceleration in mining --- consistent with economic incentives driving entry.}
\end{frame}

\begin{frame}{NDVI and EVI vs Distance to Nearest Mine}
  \centering
  \includegraphics[height=0.70\textheight]{d5a_ndvi_evi_distance_binscatter.png}

  \srcnote{Binned scatter of pixel-level NDVI (top) and EVI (bottom) vs
  distance to nearest artisanal mine, 2019. Each point = mean of a 0.5 km bin;
  point size $\propto$ pixel count; LOESS fit weighted by bin size. Landsat
  250 m composites (6.25 ha/pixel).}
\end{frame}

\begin{frame}{Agricultural Productivity Around Mine Onset --- Event Study}
  \centering
  \includegraphics[height=0.68\textheight]{d5b_ndvi_evi_event_study.png}

  \srcnote{Mean NDVI (top) and EVI (bottom) by distance band, normalised to
  $t=-1$, aligned to each location's nearest mine onset year (Barenblitt TS,
  2007--2017). Near bands = treated; far band (5--20 km) = control. A drop in
  near bands at/after $t=0$ with a flat control band is consistent with
  mine-driven vegetation loss.}
\end{frame}

\begin{frame}{Next Steps}
  \begin{columns}[T]
    \begin{column}{0.5\textwidth}
      \textbf{Empirical gaps --- data needed}
      \begin{itemize}\small\setlength\itemsep{0.35em}
        \item \textbf{Climate shocks:} does rainfall drought push people into
              mining? Needs CHIRPS annual rainfall
        \item \textbf{Formal-mine trigger:} does galamsey accelerate after a
              licence appears nearby? Needs Minerals Commission licence series
        \item \textbf{Reversion:} do mine areas recover? Needs RS panel with
              annual mine \emph{presence}
        \item \textbf{Cocoa yields:} spillovers onto COCOBOD district purchases
        \item \textbf{Labour reallocation:} is agricultural employment falling
              in high-galamsey districts? Needs census microdata
      \end{itemize}
    \end{column}
    \begin{column}{0.5\textwidth}
      \textbf{Methodological improvements}
      \begin{itemize}\small\setlength\itemsep{0.35em}
        \item \textbf{River flow direction:} replace the UTM-northing
              upstream/downstream proxy with an actual flow-direction raster
              (e.g.\ MERIT Hydro)
        \item \textbf{NDVI/EVI for cropland only:} mask to cropland pixels
              (ESA CCI / FAO land cover) to sharpen the agricultural signal
        \item \textbf{Dataset construction:} build a hexagon or
              enumeration-area panel combining mining intensity, NDVI/EVI and
              census variables
      \end{itemize}
    \end{column}
  \end{columns}
\end{frame}

%======================================================================
\appendix
\section{Annex: Additional Figures}
%======================================================================

\begin{frame}{Annex: Cumulative Mining Extent --- Districts}
  \centering
  \includegraphics[height=0.72\textheight]{d1a_districts_cumulative.png}

  \srcnote{Cumulative artisanal mining land cover per district, 2005--2019.
  Plasma colour scale (sqrt transform). Barenblitt et al.\ (2021).}
\end{frame}

\begin{frame}{Annex: Artisanal vs Industrial Mining --- Districts}
  \centering
  \includegraphics[width=0.92\linewidth]{d1b_districts_artisanal_vs_industrial.png}

  \srcnote{Left: artisanal (galamsey). Right: industrial (formal). Note the
  divergent spatial footprints. Barenblitt 2019 extent.}
\end{frame}

\begin{frame}{Annex: Top Districts by Mining Extent}
  \centering
  \includegraphics[height=0.72\textheight]{d1d_districts_extent_top20.png}

  \srcnote{Top 20 districts by cumulative artisanal mining extent (ha),
  2005--2019. A small number of districts account for the large majority of
  total galamsey.}
\end{frame}

\begin{frame}{Annex: District Trajectories Over Time}
  \centering
  \includegraphics[height=0.72\textheight]{d1c_districts_ts_top12.png}

  \srcnote{Annual new conversion (ha) for the 12 districts with highest
  cumulative extent. Onset years and growth rates vary considerably.}
\end{frame}

%======================================================================
\begin{frame}{References}
  \footnotesize
  \begin{itemize}\setlength\itemsep{0.6em}
    \item Barenblitt, A., Payton, A., Lagomasino, D., Fatoyinbo, L., et al.\
          (2021). The large footprint of small-scale artisanal gold mining in
          Ghana. \emph{Science of the Total Environment}, 781, 146644.
    \item Girard, V., et al.\ (2022). Gold-suitable geology of Africa
          (CGMW-BRGM 1:10M Geological Map derivative).
  \end{itemize}
\end{frame}

\end{document}
